\begin{theorembox}{Theorem}
\begin{theorem}[Existence of hyperbolic sector points]
\label{thm:hyperbolic-sector-existence}
For every $t\geq 0$, there exists a point $P=(u,v)$ on the upper right branch
of $x^2-y^2=1$ with
\[
  2\operatorname{HArea}(O,V,P)=t.
\]
\end{theorem}
\end{theorembox}
\begin{remark*}[Standard quantified statement]
For every admissible input satisfying the hypotheses of the statement, the
displayed definition or assertion holds in the stated domain.
\end{remark*}

\begin{remark*}[Interpretation]
The statement records the named trigonometric object or identity together with
its intended domain of use.
\end{remark*}

\begin{remark*}[Proof]
Parametrize the upper right branch by the $x$-coordinate:
\[
  P_x=(x,\sqrt{x^2-1}),\qquad x\geq 1.
\]
Let
\[
  F(x)=2\operatorname{HArea}(O,V,P_x).
\]
The sector area varies continuously with $x$: the boundary point depends
continuously on $x$, the function $\sqrt{x^2-1}$ is continuous on
$[1,\infty)$, and the area swept over a compact interval changes continuously
with the endpoint. Also $F(1)=0$, since the sector degenerates at $V$.
As $x\to\infty$, the swept sector area tends to infinity. Moreover, if
$1\leq x_1<x_2$, the sector from $V$ to $P_{x_1}$ is properly contained in
the sector from $V$ to $P_{x_2}$, so $F(x_1)<F(x_2)$. Thus $F$ is continuous,
strictly increasing, begins at $0$, and is unbounded. By the intermediate
value property, every $t\geq 0$ is equal to $F(x)$ for some $x\geq 1$.
\end{remark*}
\begin{dependencies}
\begin{itemize}
  \item \hyperref[def:hyperbolic-sector]{Hyperbolic sector}
\end{itemize}
\end{dependencies}

\begin{theorembox}{Theorem}
\begin{theorem}[Uniqueness of hyperbolic sector points]
\label{thm:hyperbolic-sector-uniqueness}
For every $t\geq 0$, the point $P$ in
\hyperref[thm:hyperbolic-sector-existence]{the existence theorem} is unique.
\end{theorem}
\end{theorembox}
\begin{remark*}[Standard quantified statement]
For every admissible input satisfying the hypotheses of the statement, the
displayed definition or assertion holds in the stated domain.
\end{remark*}

\begin{remark*}[Interpretation]
The statement records the named trigonometric object or identity together with
its intended domain of use.
\end{remark*}

\begin{remark*}[Proof]
The function $F(x)=2\operatorname{HArea}(O,V,P_x)$ is strictly increasing.
Hence $F(x_1)=F(x_2)$ implies $x_1=x_2$. Since the upper right branch point
$P_x=(x,\sqrt{x^2-1})$ is determined by $x$, the point $P$ is unique.
\end{remark*}
\begin{dependencies}
\begin{itemize}
  \item \hyperref[thm:hyperbolic-sector-existence]{Existence of hyperbolic sector points}
\end{itemize}
\end{dependencies}

\begin{corollarybox}{Corollary}
\begin{corollary}[Well-definedness of hyperbolic sine and cosine]
\label{cor:hyperbolic-sine-cosine-well-defined}
The assignments $\cosh,\sinh:\mathbb{R}\to\mathbb{R}$ in
\hyperref[def:hyperbolic-cosine]{the definition of hyperbolic cosine} and
\hyperref[def:hyperbolic-sine]{the definition of hyperbolic sine} are
well-defined functions.
\end{corollary}
\end{corollarybox}
\begin{remark*}[Standard quantified statement]
For every admissible input satisfying the hypotheses of the statement, the
displayed definition or assertion holds in the stated domain.
\end{remark*}

\begin{remark*}[Interpretation]
The statement records the named trigonometric object or identity together with
its intended domain of use.
\end{remark*}

\begin{remark*}[Proof]
Existence supplies a sector point for each $t\geq 0$, and uniqueness makes its
coordinates independent of choice. The even/odd extension then gives exactly
one value for each negative parameter.
\end{remark*}
\begin{dependencies}
\begin{itemize}
  \item \hyperref[thm:hyperbolic-sector-existence]{Existence of hyperbolic sector points}
  \item \hyperref[thm:hyperbolic-sector-uniqueness]{Uniqueness of hyperbolic sector points}
\end{itemize}
\end{dependencies}

\begin{theorembox}{Theorem}
\begin{theorem}[Fundamental hyperbolic identity]
\label{thm:fundamental-hyperbolic-identity}
For every $t\in\mathbb{R}$,
\[
  \cosh^2 t-\sinh^2 t=1.
\]
\end{theorem}
\end{theorembox}
\begin{remark*}[Standard quantified statement]
For every admissible input satisfying the hypotheses of the statement, the
displayed definition or assertion holds in the stated domain.
\end{remark*}

\begin{remark*}[Proof]
For $t\geq 0$, the point $P_t=(\cosh t,\sinh t)$ was chosen on the branch
$x^2-y^2=1$, so substituting its coordinates gives
$\cosh^2 t-\sinh^2 t=1$. For $t<0$, use the extension
$\cosh(-t)=\cosh t$ and $\sinh(-t)=-\sinh t$ to reduce to the nonnegative
case.
\end{remark*}
\begin{remark*}[Interpretation]
This identity justifies the word \emph{hyperbolic}: it follows from membership
on the unit hyperbola exactly as $\cos^2\theta+\sin^2\theta=1$ follows from
membership on the unit circle.
\end{remark*}
\begin{dependencies}
\begin{itemize}
  \item \hyperref[def:hyperbolic-sine]{Hyperbolic sine}
  \item \hyperref[def:hyperbolic-cosine]{Hyperbolic cosine}
\end{itemize}
\end{dependencies}
